% Generated draft: requires mathematical and source review.
\documentclass[11pt]{article}
\usepackage[margin=25mm]{geometry}
\usepackage{fontspec}
\setmainfont{Latin Modern Roman}
\usepackage{amsmath,amssymb,mathtools,unicode-math}
\setmathfont{Latin Modern Math}
\usepackage{enumitem,needspace,xurl,hyperref}
\setlength{\parindent}{0pt}
\setlength{\parskip}{5pt}
\setlength{\emergencystretch}{3em}
\newcommand{\N}{\mathbb{N}}
\newcommand{\Z}{\mathbb{Z}}
\newcommand{\Q}{\mathbb{Q}}
\newcommand{\R}{\mathbb{R}}
\newcommand{\C}{\mathbb{C}}
\newcommand{\F}{\mathbb{F}}
\newcommand{\A}{\mathbb{A}}
\newcommand{\PP}{\mathbb{P}}
\newcommand{\E}{\mathbb{E}}
\newcommand{\eps}{\varepsilon}
\newcommand{\dd}{\,\mathrm{d}}
\begin{document}
\section*{UnsolvedMath ID: 30003802}
\textbf{Source problem code:} catalogue entry
\subsection{Definitions and mathematical statement}
The supplied record poses a classification question, rather than a conjecture with a specified conclusion. Its original target is: describe the behavior over rings of non-unital composition algebras, including octonion isotopes. The extracted formulation additionally asks for a classification.

To explain the terminology conditionally, suppose that the base ring \(R\) is commutative with identity. An \(R\)-algebra without an assumed identity consists of an \(R\)-module \(A\) and an \(R\)-bilinear multiplication \(A\times A\to A\), written \((x,y)\mapsto xy\); associativity is not assumed. A quadratic form on \(A\) is a map \(q\colon A\to R\) satisfying \(q(rx)=r^2q(x)\) for all \(r\in R\), \(x\in A\), such that
\[
 b_q(x,y)=q(x+y)-q(x)-q(y)
\]
is \(R\)-bilinear. The composition identity is
\[
 q(xy)=q(x)q(y)\qquad\text{for all }x,y\in A.
\]
These identities explain a possible framework; they do not specify the intended class of composition algebras. The source does not state the required finiteness, projectivity, rank, or nondegeneracy conditions, nor does it explicitly restrict the base rings to commutative rings with identity. Thus these conditional definitions must not be read as additional hypotheses supplied by the problem.

In a usual algebraic sense, an isotopy between two \(R\)-algebras \(A\) and \(B\) is a triple of \(R\)-linear isomorphisms \(f,g,h\colon A\to B\) satisfying
\[
 h(xy)=f(x)\mathbin{\cdot_B}g(y)\qquad\text{for all }x,y\in A.
\]
An octonion isotope is an algebra isotopic to an octonion algebra. The record does not specify whether its intended isotopies must satisfy additional conditions involving quadratic forms, or which definition of an octonion algebra over a ring is intended.

\textbf{Statement of the research target.} Describe and classify non-unital composition algebras over the intended base rings, including those arising as octonion isotopes. A precise quantified classification assertion cannot be recovered from the supplied record: the admissible rings and algebraic data, the equivalence relation for classification, and the requested invariants or structural conclusions are unspecified. In particular, ``non-unital'' could mean that an identity is not part of the data or that no identity exists; these interpretations give different classes.
\subsection{Short English statement}
The problem asks for a description and classification of non-unital composition algebras over rings, including octonion isotopes. The supplied statement leaves the base rings, norm conditions, meaning of non-unitality, and classification equivalence unspecified, so it determines a research direction but not a unique precise conjecture.
\subsection{Sources}
\begin{itemize}
\item Supplied local UnsolvedMath registry record, permanent ID \(30003802\), title \emph{Non-Unital Composition Algebras over Rings}, source problem code \texttt{catalogue entry}, in the Oberwolfach Reports Open Problems collection. No individual proposer is supplied; the proposed year is recorded as 2018.
\item Supplied report citation: \emph{Quadratic Forms and Related Structures over Fields (2019)}, Oberwolfach Reports, pp.~1235--1291, DOI \texttt{10.4171/owr/2018/21}. Supplied source link: \url{https://doi.org/10.4171/owr/2018/21}. The background attributes the problem to workshop open-problem discussions documented in Oberwolfach Reports. The citation's 2019 date and the record's 2018 year are retained without reconciliation.
\item Supplied literature reference: \emph{Composition algebras over rings}, \url{https://doi.org/10.1017/9781009426862.005}.
\item Supplied literature reference: \emph{Oberwolfach Report 21}, \url{https://doi.org/10.4171/owr/2018/21}.
\item Literature assessment dated exactly \texttt{2026-08-21}: the supplied curation labels the problem open, says that recent work develops composition algebras over commutative rings, and reports no general resolution of the extracted question about nonunital composition algebras. This is supplied dated curation, not an independently verified literature assessment; its background explicitly requests independent verification of time-sensitive status and cited claims.
\item The registry gives the statement status as \texttt{corrected\_verified} and reports that the clean statement was checked against the official report and preserves its hypotheses, notation, and open target. This is a supplied verification claim, not verification performed for this draft. No source page locating the particular problem within the report is supplied.
\end{itemize}
\end{document}
